\scanpage{8}
\[
\begin{aligned}
&\operatorname{Tr}vuvu^{-1}=0\text{ ou }1\\
-&\operatorname{Tr}v\,u(v)=\det(v-u(\unsure{w|v}))\in\{0,1\}
\end{aligned}
\]

\[
\operatorname{Tr}(uv)\in\{0,1\}\qquad\forall u,v\in S
\]

\margin{\uline{NB} $\operatorname{Tr}u^2=1$, $\operatorname{Tr}uu'=0$}

OK si $v=u$ ou $v=u'$\qquad d'autre part,

\[
\begin{aligned}\operatorname{Tr}uv'=\operatorname{Tr}u'v&=\\[-2pt]&\operatorname{Tr}uv+1\\&\ne\operatorname{Tr}uv\end{aligned}
\]

Pour les dix autres, ceux-ci forment un ens. à $10$ éléments sur lequel $k[u]^*/k^*$ opère or $k[u]^*$ est cyclique d'ordre $4^2-1=15$, $k^*$ est ss groupe d'ordre $3$, donc $k[u]^*/k^*$ est cyclique d'ordre $5$. On vérifie que $u^5=\alpha.1$ donc $u$ est d'ordre $15$ ($\alpha$ étant d'ordre $3$) donc $u$ \deleted{est}\inserted{définit} un générateur de $k[u]^*/K^*$.

Ceci dit, pts fixes de \deleted{\unsure{Le commutant}}\inserted{$\operatorname{End}(V)=M$.} \deleted{\unsure{dans}} \inserted{l'un des} de $k[\underline u]^*/k^*$, i.e. \unsure{le|les} commutant de $k[u]$ dans $M$, \unsure{du|de} $u$ est réduit à $k[u]$. On a vu que $k[u]\cap S=\{u,u'\}$ donc $u$ ne fixe aucun pt de $S\smallsetminus\{u,u'\}$, qui se décompose donc en deux orbites d'ordre $5$. Donc $\operatorname{Tr}(uv)$ pour $v\in S\smallsetminus\{u,u'\}$ ne peut prendre que deux valeurs au plus (une par orbite)

\margin{\uline{Cor 1} $M^*/k^*\unsure{=|\simeq|=:}G$ transitif sur $S$\\\uline{Cor. 2.} $G$ simplement transitif}

% archive sheet number: 7
